Sparse \method{} reduces the left/right regional objectives by $50.87\%$ /\allowbreak $50.09\%$, compared with $38.15\%$ /\allowbreak $29.99\%$ for deterministic hottest-point placement. The primary multigrid topological selector gives $53.30\%$ /\allowbreak $52.15\%$. Averaged across both cases, complete sparse \method{} design takes $3.66$\,s, versus $8.38$\,s for the primary multigrid setting. The fastest tested multigrid setting (57 nodes along the first axis) takes $2.51$\,s and gives $53.50\%$ /\allowbreak $52.53\%$ reduction. Retaining the outer mesh in one process per complete adaptive design reduces wall time to $1.22$\,s for sparse \method{}, $1.41$\,s for coarse multigrid, and $7.13$\,s for primary multigrid. All 70 repeated designs have exactly the same six sequential cavity centers as their original runs, permitting reuse of the independent objective evaluations. These results demonstrate an advantage over hottest-point placement; the full budget sweep in \cref{fig:regional,tab:regional} qualifies comparisons with volumetric topological derivatives.
